\section{Common-task calibration and causal simulation}
\label{sec:composition}
A causal experiment morphism consists of a parameter-independent report transformation, strategy lift, internal simulator state, compatible actions, physical clock and task readout. Its deficiency is measured on the complete stopped and scored law, including the observer's fresh coins. If that total variation defect is at most $\eps$ and $0\le\ell\le L_\ell$, then
\begin{equation}\label{eq:morphism}
 R_{\rm source}(\rho(B))\le R_{\rm target}(B)+L_\ell\eps.
\end{equation}
Indeed run the lifted observer and simulator jointly, compare their bounded common-task expectations, and take an infimum. Two such defects add by contraction and the triangle inequality. If the simulator uses $R$ internal states and the observer uses $M$ labels, their direct product uses at most $RM$ labels. Mean resource transfer requires its own pathwise or potential certificate; it does not follow from total variation. Theorem~\ref{thm:perturbation} supplies one such certificate with protocol modes that need not stop at each call. A one-way morphism does not imply a converse risk inequality.

If the final prediction differs in root mean square by at most $\delta_r$ under the same law, Minkowski's inequality gives
\[
 \sqrt{\E\|Z-\widehat r\|^2}\le\sqrt{\E\|Z-r\|^2}+\delta_r.
\]
Root errors add before squaring. An upper tolerance $\delta_r$ is not an information-theoretic lower bound. The following attained-error experiment makes the distinction explicit.

\begin{theorem}[An attained joint resource--risk law]\label{thm:joint}
Use either gate realization with $B_0=0$ and fixed calibration/geometric constants. Before the gate protocol, one paid call prepares a fair sign $U\in\{-1,1\}$ and reports it unless it is erased with probability $2\eps$, $0\le\eps\le1/2$. The core starts independently of $U$; none of its kernels or audit coordinates reads $U$. At the final paid audit score the vector $(Z,U,\xi\delta)$, where $\xi\in\{-1,1\}$ is a fixed unknown calibration alternative with identical pre-audit laws. The sign response is restricted to $[-1,1]$. Count the early call and final audit, so a total budget $N$ leaves $N-2$ expected raw calls for the core.

For $N\ge\kappa(p)+2$ and $M\ge13$, put
\[
 m_M=\left\lfloor\frac{M-1}{3}\right\rfloor-2\ge2,\qquad
 b=\frac{N-2}{\kappa(p)}.
\]
The worst-model autonomous risk for this same task satisfies
\begin{equation}\label{eq:joint-law}
 \mathcal R(N,M,\delta,\eps)
 \asymp F_{d,\alpha}(b,m_M)+\delta^2+\eps.
\end{equation}
The constants are independent of $p,N,M,\delta,\eps$ for bounded $\delta$ and fixed core certificates. The early channel's causal deficiency to an exact immediate sign-report interface is exactly $\eps$. It must complete that virtual report before entering the core.
\end{theorem}
\begin{proof}
The core lower holds after every allowed history, including knowledge of an unerased $U$ and all previous coins. Additional independent randomness cannot defeat its actual-occupation bound. Projecting all terminal values of an $M$-label program to the core gives at most $M-1$ core values. Consequently its core risk is at least a constant times $F_{d,\alpha}(b,M-1)$. Conditional on erasure, $U$ remains fair and independent of all permitted reports and actions: this follows inductively because the core kernels do not read $U$. Its unavoidable squared loss is then one, contributing $2\eps$. The two calibration alternatives have identical pre-audit laws; averaging their risks for any response $r$ gives $\E r^2+\delta^2$. Taking the worse alternative yields at least $\delta^2$. These lower terms can be combined on the same task and model; alternatively their maximum bounds a third of their sum.

For the upper use one initial label, three core lanes for $U=+1,-1,$ or erased, and in each lane two live protocol labels and $m_M$ terminal labels. The total is $1+6+3m_M\le M$. Run the common core program in each lane, predict the reported sign or zero on erasure, and predict zero for the calibration coordinate. Its risk is at most a constant times $F_{d,\alpha}(b,m_M)+2\eps+\delta^2$, and its expected total calls are at most $N$. Since $m_M\ge2$ and $M-1<3m_M+9\le8m_M$, the two core widths are comparable uniformly, proving \eqref{eq:joint-law}.

On erasure any simulator completing the sign report immediately has no information about $U$ and errs with probability at least $1/2$. Completion by an independent fair sign attains that error. The complete joint law of true and virtual signs is therefore at total variation distance exactly $(2\eps)/2=\eps$ from the exact-sign law. Later physical audits cannot be used to complete an earlier virtual report. This proves attained deficiency, not a lower bound inferred from an arbitrary allowed tolerance.
\end{proof}

The total budget in \eqref{eq:joint-law} concerns a stopped protocol, not a deterministic horizon. Calibration and numerical errors are distinguished: $\delta$ in this theorem is an actual hidden alternative, while numerical routing errors require the separate feasible-cost certificates of Section~\ref{sec:perturbation}. Deterministic sample caps, unbounded losses, large-deviation equivalence and infinite-horizon limit interchange are not consequences of the displayed total variation inequality.
